%==============================================================================
% Capitulo 11: Consideraciones eticas y aspectos legales
%==============================================================================
\section{Consideraciones eticas y aspectos legales}

\subsection{Marco normativo aplicable}

El sistema desarrollado se enmarca normativamente en la Ley~25.326 de Proteccion de los Datos Personales de la Republica Argentina \parencite{Ley25326_2000}, en su decreto reglamentario y en las disposiciones complementarias emitidas por la Agencia de Acceso a la Informacion Publica en su caracter de autoridad de aplicacion. Dado que el sistema realiza transferencia de fragmentos textuales hacia infraestructura de inferencia operada por proveedores con sede en los Estados Unidos, resultan aplicables las consideraciones del articulo 12 de la mencionada ley sobre transferencia internacional de datos personales y, de manera complementaria, los lineamientos del Reglamento General de Proteccion de Datos de la Union Europea \parencite{GDPR2016} en aquellos casos que pudieran involucrar a residentes europeos. Argentina mantiene un regimen de adecuacion reconocido por la Comision Europea, lo cual facilita los flujos transfronterizos bidireccionales bajo determinadas condiciones.

\subsection{Principios aplicados al tratamiento de datos}

El diseno del sistema se atiene de manera explicita a cinco principios fundamentales del tratamiento de datos personales. El principio de licitud establece que todo tratamiento debe ampararse en una base juridica valida; en este caso, dicha base se encuentra en el legitimo interes de la organizacion para la gestion interna de los incidentes informaticos reportados por sus propios usuarios en ejercicio de sus funciones laborales. El principio de finalidad determinada limita el uso de los datos a la finalidad declarada de registro y derivacion de incidentes, prohibiendo cualquier uso secundario no compatible. El principio de minimizacion exige recolectar unicamente los datos estrictamente necesarios, razon por la cual el sistema no almacena identificadores personales mas alla del nombre de usuario corporativo y descarta cualquier dato adicional capturado incidentalmente. El principio de exactitud se materializa mediante la posibilidad ofrecida al usuario de revisar y corregir el contenido del incidente antes de su persistencia definitiva. El principio de limitacion del plazo de conservacion se traduce en una politica de retencion de 90 dias para los registros operativos y de un ano para los datos de incidentes resueltos.

\subsection{Transferencia internacional y pseudonimizacion}

La utilizacion del modelo Gemini~2.5 Flash implica la transmision de fragmentos textuales descriptivos de los incidentes hacia infraestructura del proveedor con sede en los Estados Unidos. Para mitigar el riesgo asociado a esta transferencia, el sistema aplica un procedimiento de pseudonimizacion previa a la transmision, mediante el cual se reemplazan automaticamente los nombres propios, las direcciones de correo electronico, los numeros telefonicos, los identificadores corporativos y los nombres de hosts internos por etiquetas genericas tales como \texttt{[PERSONA]}, \texttt{[EMAIL]} o \texttt{[HOST]}. Esta pseudonimizacion se realiza dentro del modulo Python desplegado en infraestructura local de la organizacion y se valida mediante un conjunto de expresiones regulares especificas del dominio, acompanadas de pruebas unitarias dedicadas. Adicionalmente, la organizacion ha suscripto el acuerdo de procesamiento de datos provisto por el proveedor del servicio de inferencia.

\subsection{Seguridad tecnica}

Las medidas de seguridad tecnica implementadas se sustentan en el modelo de la triada confidencialidad, integridad y disponibilidad (CIA) descrito por \textcite{Stallings2017_computer_security}. La confidencialidad se garantiza mediante cifrado en transito sobre TLS~1.3 entre todos los componentes del sistema y mediante cifrado en reposo de la base de datos PostgreSQL utilizando la extension pgcrypto para los campos sensibles. La integridad se asegura mediante el uso de identificadores firmados con HMAC-SHA-256, la validacion estricta de esquemas mediante Pydantic en cada llamada a la interfaz REST, y un registro de auditoria de toda modificacion que conserva el actor, la marca temporal y el delta del cambio. La disponibilidad se respalda mediante respaldos diarios automaticos con retencion escalonada, despliegue en contenedores con politicas de reinicio automatico ante fallas, y monitoreo de salud activo mediante endpoints de chequeo. Las credenciales de servicios externos se gestionan exclusivamente mediante variables de entorno inyectadas por el orquestador de contenedores y nunca se almacenan en el repositorio de codigo fuente.

\subsection{Consentimiento, derechos del usuario y supervision humana}

Los usuarios de la organizacion son informados, mediante una politica de privacidad accesible desde los canales de entrada, sobre las caracteristicas del sistema, los datos que se recolectan, la transferencia internacional aplicable y los derechos que les asisten conforme a la normativa vigente. Se garantiza el ejercicio de los derechos de acceso, rectificacion, actualizacion y supresion (derechos ARCO) mediante un procedimiento documentado a cargo del responsable de proteccion de datos designado por la organizacion.

En lo que respecta al consentimiento informado para los fines de investigacion academica, los 120 empleados de la organizacion fueron notificados mediante comunicacion interna formal, con descripcion explicita del proposito del estudio, del tipo de datos tratados, del caracter pseudonimizado de los registros y del derecho a requerir la exclusion de sus incidentes del corpus de validacion. El procedimiento fue supervisado por la direccion academica del trabajo y avalado institucionalmente mediante una carta de autorizacion emitida por la organizacion adoptante.

En relacion con el protocolo de observacion naturalista aplicado durante la fase experimental ---en el cual los operadores no fueron informados en tiempo real de que sus tiempos de registro estaban siendo medidos, con el objeto de eliminar el efecto Hawthorne---, corresponde dejar constancia de que, una vez concluida la recoleccion de datos, se realizo una sesion de esclarecimiento posterior (\textit{debriefing}) con los dos operadores participantes. En dicha sesion se explico el proposito del estudio, se detallo el procedimiento de medicion empleado y se ofrecio a cada operador la posibilidad de solicitar la exclusion de sus registros individuales del corpus de validacion. Ninguno de los operadores ejercio este derecho. Este procedimiento se alinea con las recomendaciones de la Asociacion Americana de Psicologia (APA, 2017) para investigaciones que emplean observacion sin conocimiento previo de los participantes, y fue supervisado por el director del trabajo.

En concordancia con las recomendaciones contemporaneas sobre sistemas que incorporan inteligencia artificial, se preserva en todo momento la posibilidad de supervision humana significativa. Ninguna decision que afecte derechos sustanciales de los usuarios se toma de manera completamente automatizada sin posibilidad de revision, y el sector responsable conserva la potestad de modificar la clasificacion inicial cuando lo considere apropiado, registrandose tal modificacion en la tabla de trazabilidad descrita en el Capitulo~5. Este mecanismo materializa el principio de la persona en el bucle (\emph{human in the loop}) como salvaguarda frente a errores potenciales del clasificador automatico.

\newpage
